\chapter{Superseded AVX-led Introduction}
\label{ch:introduction}

Thirty-two bytes sit in one machine register. Give them names from left to
right:

\[
a_0,a_1,\ldots,a_{30},a_{31}.
\]

We want the order reversed. The last byte must come first, the first must come
last, and every byte between them must cross to its matching place. One vector
instruction can reverse the bytes inside each sixteen-byte half. A second can
exchange the two halves. In the restricted single-register language studied
later, two instructions do the job.

There is a small pleasure in watching thirty-two separate positions fall into
order after two machine steps. But the program leaves a more difficult
question behind. Two instructions are enough. Are two necessary?

The difference between those sentences is the subject of this book. The
first points to something we can hold: a sequence of instructions. The second
reaches over every cheaper sequence we might have written and says that none
works. A tiny program has opened a universal question.

A compiler cannot settle it merely by emitting two instructions. Its rules
may have found them without ruling out one. Tests cannot settle it either.
They can expose a bad reversal, but even a great many passing inputs leave
other inputs untouched. Inspection becomes least dependable where vector
instructions become most powerful. A single instruction can select, blend,
shift, and reorder more data than its printed name suggests.

We need a way to turn \emph{no cheaper program works} into an object a reader
can examine. Before we build that object, we must say exactly what the claim
means.

\section{What shortest means}
\label{sec:intro-shortest}

An instruction takes a machine state and returns another machine state. In a
small problem, that state may be one register. A program applies one
instruction after another. It is therefore a composition of functions: the
output of one function becomes the input of the next.

This description comes before the traditional RISC/CISC division. Those names
gather several design tendencies, but they do not supply the semantics of an
instruction. We will instead name each feature that changes a proof. Are the
operands explicit? May an instruction read memory? Does it change flags or
return several results? Which encodings belong to the program language? A
RISC-V instruction and an x86-64 instruction can then be
compared as exact state transformations, not as mascots for two camps.

Take four tagged bytes \(a,b,c,d\). A reversal must return \(d,c,b,a\). The
tags make the motion visible, but this row is only an example. Two programs
are \keyterm{equivalent} when they return the same result for every input
allowed by the problem. One matching row can suggest equivalence. It cannot
establish it.

The word \emph{shortest} needs another choice. Which programs are allowed to
compete? May an instruction read memory? May it use a second source register?
Are all immediate constants available? Does every instruction cost one, or
do we measure latency on a named processor? The \keyterm{program language}
sets the allowed instructions and operand forms. The \keyterm{cost model}
says how their costs are compared. Together they determine the
\keyterm{scope} of the claim.

Scope is part of the proposition. A two-instruction sequence may be shortest
among one-source register operations and lose when a memory lookup joins the
language. It may minimize instruction count and lose under a latency model.
Nothing paradoxical has happened. We asked a different question.

Once the scope is fixed, the claim that no cheaper program works is called
\keyterm{minimality}. It splits into two bounds. A program of cost \(k\) is a
witness that the best cost is at most \(k\). To show that the best cost is at
least \(k\), we must rule out every equivalent
program of lower cost. The witness descends from above. The lower bound rises
from below. When they meet, the optimum is pinned to one number.

This meeting is simple to draw and hard to earn. Finding a good program and
showing that no better one exists are different mathematical acts.

\section{Search and the missing receipt}
\label{sec:intro-history}

Programmers have pursued the smallest machine programs for decades. In 1987,
Henry Massalin described a superoptimizer, a system that searched an
instruction set for the shortest program computing a given function
\autocite{massalin1987}. Sorav Bansal and Alex Aiken later used brute-force
search to generate large collections of peephole rewrites automatically
\autocite{bansalaiken2006}. This kind of search is called
\keyterm{superoptimization}. Recent systems still press against the same
boundary. HieraSynth divides the search space without discarding candidates;
when an exhaustive run finishes, it supports an optimality claim within its
stated cost model \autocite{lubodik2025}.

The history matters here for one reason. Searching for short programs is not
new, and neither is the desire to know that a search has reached the bottom.
What changes from system to system is the warrant. A fast search may find a
remarkable program without supplying a lower bound. An exhaustive search can
support a lower bound only for the language its encoding actually covers. A
complete search can still optimize a cost model that poorly matches the
machine we care about.

A further gap often remains after a correct search finishes. The search may
leave no compact object for an independent program to check. A log saying
\emph{no candidate found} asks us to trust the searcher, its pruning rules,
and its implementation. Running the search again may be worthwhile. It is not
the same as checking a receipt from the search.

This distinction contains a beautiful asymmetry. One short sequence can show
that a program exists. Nonexistence must somehow account for an entire space
of alternatives. The proof of absence may be much larger than the witness,
but if it has the right form, a simpler program can still check it.

The aim of this book is narrower than inventing optimization and more exact
than reporting successful searches. For each machine claim, we want to know
what bytes carry its evidence, what program checks those bytes, how that
program behaves when the evidence is wrong, and where the claim stops.

\section{The unit of belief}
\label{sec:intro-object}

An \keyterm{artifact} is the stored evidence for a claim. It may be an
instruction sequence, a formula, a refutation, or the output of a finite
enumeration. A \keyterm{checker} reads that evidence and gives a result that
depends on whether the evidence supports the stated claim. Alter the witness
or truncate the proof, and the checker must object.

That requirement also needs evidence. A \keyterm{negative control} is a
deliberately false claim or damaged artifact sent through the same route. The
control should fail. If it passes, success on the original object carries
little force. A checker that agrees with everything is agreeable company and
poor evidence.

This book binds the pieces in an \keyterm{object}. An object records the
claim, its scope, its evidence category, its artifact, and the expected
success and failure conditions. The \keyterm{ledger} lists those objects.
Chapters take status and scope from it, so a sentence cannot quietly turn a
result about bytes into a statement about every width.

Return once more to the reversal. Its two-instruction sequence is the
artifact for the upper bound. A checker can replay the sequence under the
defined instruction functions. The lower bound needs another artifact, one
that rules out every one-instruction program in the stated language. A false
one-instruction claim or damaged proof tests whether that route can lose. The
artifact, checker, scope, object, ledger entry, and control are not six pieces
of bureaucracy. They are six answers to the question, \emph{Why should I
believe this particular sentence?}

The ledger does not teach the reason. Every load-bearing theorem therefore
has two proofs with different work to do. A reader proof exposes the idea in
a small case, invariant, contradiction, or replay. Its machine counterpart
checks the full claim in its named scope. One gives understanding without full
coverage. The other can give coverage without understanding. We need both.

Not every useful result is a theorem. A finite search may produce a checked
table. An unfinished argument may isolate an exact missing obligation. The
ledger keeps these objects in their proper categories. A smooth narrative is
not permission to strengthen the evidence.

Evidence also has depth. A solver may report that a negated claim has no
solution. It may export a certificate for another program to check. A proof
assistant may admit a term only after a small kernel checks every step. These
routes place different programs inside the boundary of trust. We will name
the route we have rather than borrow the force of one we do not.

Here the machinery returns us to the original wonder. Thirty-two bytes and
two instructions were easy to write down. The phrase \emph{no one-instruction
program} was invisible until we gave its universe a boundary and its evidence
a form. Exactness has made a much larger object available to thought.

\section{How the argument moves}
\label{sec:intro-reading}

Most chapters begin with a machine problem that fits in the hands: a byte
permutation, a population count, a bitwise identity, or a short schedule. We
work a small instance before naming its general structure. Then we state the
scoped claim, show the reader argument, inspect the artifact, and try to break
the checker. The boundary of one result supplies the next question.

The first part stays below any particular ISA. It constructs words, state,
memory, operand forms, and straight-line programs. Concrete RISC-V and x86-64
examples enter only after those common objects are available. Their
differences will test the definitions rather than determine them.

The exercises continue the argument. Some reproduce a result. Some damage an
artifact on purpose. Some change a width, instruction family, or cost model
until an earlier proof no longer applies. The harder exercises ask for the
next object the book does not yet possess.

Before two instruction sequences can be equivalent, their instructions must
act on something exact. Before we can rule out a one-instruction program, we
need a precise input space and precise operations. The register at the start
therefore sends us inward, from thirty-two bytes to one byte, and then to the
smallest things the machine can distinguish: a fixed number of zeros and
ones.
